% Shared preamble for ECON 5280 lecture notes.
\usepackage[T1]{fontenc}
\usepackage[utf8]{inputenc}
\usepackage[sc]{mathpazo}
\linespread{1.04}
\usepackage{microtype}

\usepackage[a4paper,margin=1in,headheight=15pt]{geometry}
\usepackage{amsmath,amssymb,amsthm,mathtools,bm,mathrsfs}
\usepackage{graphicx}
\graphicspath{{../assets/}}
\usepackage{booktabs,longtable,array,multirow}
\usepackage{enumitem}
\usepackage{float}
\usepackage{caption}
\usepackage{subcaption}
\usepackage[dvipsnames]{xcolor}
\usepackage{listings}
\usepackage[most]{tcolorbox}
\usepackage{titlesec}
\usepackage{fancyhdr}
\usepackage{url}
\usepackage[hidelinks]{hyperref}
\usepackage[nameinlink,noabbrev]{cleveref}

\definecolor{econblue}{HTML}{17365D}
\definecolor{econlight}{HTML}{F1F5F9}
\definecolor{econgreen}{HTML}{285943}
\definecolor{codebg}{HTML}{F7F7F5}
\definecolor{codecomment}{HTML}{4F6F52}
\definecolor{codestring}{HTML}{8A3B12}

\hypersetup{
  colorlinks=true,
  linkcolor=econblue,
  citecolor=econgreen,
  urlcolor=econblue,
  pdftitle={ECON 5280: Causal Inference and Machine Learning},
  pdfauthor={Junlong Feng}
}

\setlist{itemsep=0.25em,topsep=0.5em}
\setcounter{secnumdepth}{3}
\setcounter{tocdepth}{2}
\allowdisplaybreaks
\emergencystretch=2em

\titleformat{\chapter}[display]
  {\normalfont\bfseries\color{econblue}}
  {\filleft\Large\chaptertitlename\ \thechapter}
  {0.6ex}
  {\titlerule\vspace{0.9ex}\Huge\raggedright\hyphenpenalty=10000\relax}
\titlespacing*{\chapter}{0pt}{-10pt}{30pt}

\pagestyle{fancy}
\fancyhf{}
\fancyhead[L]{\small\nouppercase{\leftmark}}
\fancyhead[R]{\small ECON 5280}
\fancyfoot[C]{\thepage}
\renewcommand{\headrulewidth}{0.4pt}

\newcommand{\E}{\mathbb{E}}
\newcommand{\Var}{\operatorname{Var}}
\newcommand{\Cov}{\operatorname{Cov}}
\newcommand{\Corr}{\operatorname{Corr}}
\newcommand{\Prb}{\mathbb{P}}
\newcommand{\N}{\mathcal{N}}
\newcommand{\dto}{\overset{d}{\longrightarrow}}
\newcommand{\pto}{\overset{p}{\longrightarrow}}
\newcommand{\indep}{\mathrel{\perp\!\!\!\perp}}
\newcommand{\1}{\mathbf{1}}
\newcommand{\R}{\mathbb{R}}

\theoremstyle{plain}
\newtheorem{theorem}{Theorem}[chapter]
\newtheorem{proposition}[theorem]{Proposition}
\newtheorem{lemma}[theorem]{Lemma}
\newtheorem{claim}[theorem]{Claim}
\theoremstyle{definition}
\newtheorem{definition}{Definition}[chapter]
\newtheorem{assumption}{Assumption}[chapter]
\newtheorem{example}{Example}[chapter]
\theoremstyle{remark}
\newtheorem{remark}{Remark}[chapter]

\newtcolorbox{important}{
  enhanced,
  breakable,
  colback=econlight,
  colframe=econblue,
  boxrule=0.7pt,
  arc=1.5mm,
  left=2mm,right=2mm,top=1.5mm,bottom=1.5mm,
  title=Important,
  fonttitle=\bfseries
}

\lstdefinestyle{econR}{
  language=R,
  basicstyle=\ttfamily\small,
  backgroundcolor=\color{codebg},
  frame=single,
  rulecolor=\color{black!18},
  breaklines=true,
  breakatwhitespace=false,
  columns=fullflexible,
  keepspaces=true,
  showstringspaces=false,
  commentstyle=\color{codecomment}\itshape,
  stringstyle=\color{codestring},
  keywordstyle=\color{econblue}\bfseries,
  numbers=left,
  numberstyle=\tiny\color{black!50},
  numbersep=7pt,
  xleftmargin=1.4em,
  framexleftmargin=1.1em,
  aboveskip=0.8em,
  belowskip=0.8em
}
\lstnewenvironment{rcode}[1][]
  {\lstset{style=econR,#1}}
  {}

\captionsetup{font=small,labelfont=bf}
\crefname{assumption}{assumption}{assumptions}
\Crefname{assumption}{Assumption}{Assumptions}
